% ============================================================
% §4.5 Extension to General Local Search Operators
% Generalizes the random-search-specific results in §4.3–§4.4
% to an abstract operator class covering ε-greedy and
% local perturbation search.
% ============================================================
\subsection{Extension to General Local Search Operators}
\label{sec:general-local-search}

The preceding analysis
(\cref{sec:sde-foundations,sec:scheduling-optimality}) establishes the
gain factorization $G_t(K)=\sigma_t\,a_K$ and its scheduling
consequences for the special case of pure random search.
The key structural insight, however, is that the scheduling results
depend only on three properties of the gain sequence:
monotonicity, non-accelerating marginals, and $t$-independence.
Any local search operator sharing these properties inherits
the same water-filling (\cref{prop:offline}) and
Jensen-gap/dual-threshold (\cref{prop:online}) guarantees.
This section formalizes this observation, extending the theory to
$\epsilon$-greedy exploration and local perturbation search under
four regularity conditions.

% -----------------------------------------------------------
\subsubsection{Abstract Local Search Framework}
\label{sec:abstract-ls}

A local search operator $\mathcal{L}$ at timestep~$t$ produces iterates
$\xi^{(0)},\xi^{(1)},\ldots,\xi^{(K)}$ in $\mathbb{R}^d$,
starting from an initial draw $\xi^{(0)}\sim\mathcal{N}(0,I)$.
The \emph{gain} from $K$ iterations is
\begin{equation}
  G_t^{\mathcal{L}}(K)
  \;:=\;
  \mathbb{E}\!\Bigl[
    \max_{0\le j\le K}
    \Psi_t\!\bigl(\bar{X}_{t-h}+g_t\sqrt{h}\,\xi^{(j)};\,c\bigr)
  \,\Big|\, X_t,c\Bigr]
  -\mu_t,
  \label{eq:gain-general}
\end{equation}
where $\mu_t=\Psi_t(\bar{X}_{t-h};c)$ as before.
The following four assumptions suffice for all subsequent results.

\begin{assumption}[Local search regularity]
\label{asm:local-search}
A local search operator $\mathcal{L}$ satisfies:
\begin{enumerate}
\item[\textbf{(A1)}]
  \emph{Strict monotone improvement.}\;
  $G_t^{\mathcal{L}}(K)$ is strictly increasing in~$K$ for $K\ge 1$.
\item[\textbf{(A2)}]
  \emph{Diminishing returns.}\;
  $G_t^{\mathcal{L}}(K{+}1)-G_t^{\mathcal{L}}(K)$ is bounded above
  by a non-increasing sequence, i.e., the marginal gains do not
  accelerate.
\item[\textbf{(A3)}]
  \emph{Sensitivity scaling.}\;
  In the linearized regime ($g_t\sqrt{h}\to 0$),
  \begin{equation}
    G_t^{\mathcal{L}}(K)
    \;=\; \sigma_t\,\phi_K^{\mathcal{L}}
    \;+\; O_K(g_t^2\,h),
    \label{eq:sensitivity-scaling}
  \end{equation}
  where $\phi_K^{\mathcal{L}}\ge 0$ depends on the operator
  $\mathcal{L}$ and the budget $K$ but \emph{not} on the
  timestep~$t$, and the implicit constant in $O_K(\cdot)$ may
  depend on $K$, $d$, and $L_t$.
\item[\textbf{(A4)}]
  \emph{Rotational equivariance.}\;
  For any orthogonal matrix $Q\in\mathbb{R}^{d\times d}$, the joint
  distribution of $\{Q\xi^{(j)}\}_{j\ge 0}$ under $\mathcal{L}$
  equals that of the iterates produced by $\mathcal{L}$ starting
  from $Q\xi^{(0)}$.
\end{enumerate}
\end{assumption}

\begin{remark}[Grounding of assumptions]
\label{rmk:grounding}
Each assumption captures a natural property of practical operators.
\textbf{(A1)} holds whenever the operator
maintains a running best-so-far (a ``keep best'' wrapper) and each
iteration has a positive probability of improving the incumbent.
\textbf{(A2)} reflects the general principle that the first
few iterations are most productive: for random search this follows from
order-statistic theory ($a_K$ is concave;
\citealt{david2003order}, Ch.~4); for gradient-based methods it
follows from convergence on smooth objectives; for MCMC it follows from
mixing-time arguments.
\textbf{(A3)} arises because, in the small-noise regime,
$\Psi_t(\bar{X}_{t-h}+g_t\sqrt{h}\,\xi;c)
= \mu_t+\sigma_t\langle u_t,\xi\rangle + O(g_t^2 h)$
where $u_t:=\nabla\Psi_t/\|\nabla\Psi_t\|$ is a fixed unit vector
(\cref{prop:taylor}).
Any local search on $\xi$ therefore reduces to optimizing a linear
function $f(\xi)=\langle u,\xi\rangle$, and the gain from this
optimization depends on the search method and budget but not on the
direction~$u$ or the scaling factor~$\sigma_t$, beyond a
multiplicative factor.
\textbf{(A4)} holds for all three operators below, since each uses
only isotropic primitives ($\mathcal{N}(0,I)$ draws); it ensures that the
projected gain $\mathbb{E}[\max_j\langle u,\xi^{(j)}\rangle]$ does
not depend on the direction~$u$.
\end{remark}

% -----------------------------------------------------------
\subsubsection{General Gain Factorization}
\label{sec:general-gain}

The location-scale gain decomposition of \cref{thm:loc-scale}(iii)
extends to any operator satisfying \cref{asm:local-search}.

\begin{theorem}[General gain factorization]
\label{thm:general-gain}
Under \cref{asm:smooth,asm:local-search}, for any local search operator
$\mathcal{L}$ with operator-specific gain sequence
$\{\phi_K^{\mathcal{L}}\}_{K\ge 0}$, the expected gain from $K$
iterations at timestep~$t$ satisfies
\begin{equation}
  G_t^{\mathcal{L}}(K)
  \;=\;
  \sigma_t\,\phi_K^{\mathcal{L}}
  + O_K(g_t^2\,h),
  \label{eq:general-gain-factorization}
\end{equation}
where $\sigma_t=g_t\sqrt{h}\,\|\nabla\Psi_t(\bar{X}_{t-h};c)\|$
and the implicit constant in $O_K(\cdot)$ depends on $K$, $d$,
and $L_t$ but not on $g_t$ or $\nabla\Psi_t$.
The sequence $\phi_K^{\mathcal{L}}$ is strictly increasing for
$K\ge 1$ (by~A1), has non-accelerating marginal gains (by~A2),
and satisfies $\phi_0^{\mathcal{L}}=0$.
\end{theorem}

\begin{proof}
By \cref{prop:taylor},
$\Psi_t(\bar{X}_{t-h}+g_t\sqrt{h}\,\xi;c)
=\mu_t+\sigma_t\langle u_t,\xi\rangle+R_t(\xi)$,
where $u_t=\nabla\Psi_t(\bar{X}_{t-h};c)/\|\nabla\Psi_t(\bar{X}_{t-h};c)\|$
and $|R_t(\xi)|\le\tfrac{1}{2}L_t g_t^2 h\|\xi\|^2$.
Since $\mu_t$ is constant across iterates, substituting
into~\eqref{eq:gain-general} gives
\begin{align*}
  G_t^{\mathcal{L}}(K)
  &= \mathbb{E}\Bigl[
       \max_{0\le j\le K}
       \bigl(\sigma_t\langle u_t,\xi^{(j)}\rangle + R_t(\xi^{(j)})\bigr)
     \,\Big|\,X_t,c\Bigr] \\
  &= \sigma_t\,\mathbb{E}\Bigl[
       \max_{0\le j\le K}\langle u_t,\xi^{(j)}\rangle
     \,\Big|\,X_t,c\Bigr]
     + O_K(g_t^2 h),
\end{align*}
where the second equality uses the elementary inequality
$|\max_j(a_j+b_j)-\max_j a_j|\le\max_j|b_j|$
(as in \cref{thm:loc-scale}(iii)), the non-negativity $\sigma_t\ge 0$
to factor $\sigma_t$ out of the max, and the remainder bound
$\mathbb{E}[\max_{0\le j\le K}|R_t(\xi^{(j)})|]
\le\tfrac{1}{2}L_t g_t^2 h\,\mathbb{E}[\max_{0\le j\le K}\|\xi^{(j)}\|^2]$.
The subscript~$K$ in $O_K(\cdot)$ reflects that the constant depends on
$\mathbb{E}[\max_j\|\xi^{(j)}\|^2]$, which may grow with~$K$ for
operators whose iterates accumulate drift.

By~(A4), the joint distribution of $\{\xi^{(j)}\}$ is rotationally
equivariant.
Hence, for any fixed unit vector $u_t$, the projected iterates
$\{\langle u_t,\xi^{(j)}\rangle\}$ have the same distribution as
$\{\langle e_1,\xi^{(j)}\rangle\}$ (where $e_1$ is any fixed
coordinate direction).
Therefore the expected projected gain
$\mathbb{E}[\max_{0\le j\le K}\langle u_t,\xi^{(j)}\rangle]$
depends on $\mathcal{L}$ and $K$ only; we denote it
$\phi_K^{\mathcal{L}}$.
Strict monotonicity of $\phi_K^{\mathcal{L}}$
follows from~(A1), non-accelerating marginals from~(A2),
and $\phi_0^{\mathcal{L}}=0$ since
$\mathbb{E}[\langle u,\xi^{(0)}\rangle]=0$ for
$\xi^{(0)}\sim\mathcal{N}(0,I)$.
\end{proof}

% -----------------------------------------------------------
\subsubsection{Scheduling Optimality for General Operators}
\label{sec:general-scheduling}

The proofs in \cref{sec:scheduling-optimality} use monotonicity,
non-accelerating marginals, and $t$-independence of the gain
sequence.
\cref{thm:general-gain} guarantees all three for
$\phi_K^{\mathcal{L}}$, so we can replace $a_K$ throughout.

\begin{corollary}[Offline water-filling for general operators]
\label{cor:offline-general}
Under the conditions of \cref{thm:general-gain}, with
prompt-independent~$\sigma_t$,
the allocation problem
\begin{equation}
  \max_{\{K_t\}:\,\sum_t K_t=B,\;K_t\ge 1}\;
  \sum_{t=1}^{T}\sigma_t\,\phi_{K_t}^{\mathcal{L}}
  \label{eq:alloc-general}
\end{equation}
has the same structure as \cref{prop:offline}:
\begin{enumerate}
\item[\textup{(i)}]
  The optimal allocation $\{K_t^*\}$ is simultaneously optimal for
  every instance.
\item[\textup{(ii)}]
  $K_t^*$ is non-decreasing in~$\sigma_t$.
  When $\phi_K$ is strictly concave, $K_t^*$ is strictly
  increasing (water-filling); when $\phi_K$ is linear, the
  solution concentrates all budget on $\arg\max_t\sigma_t$.
\item[\textup{(iii)}]
  Non-uniform allocation weakly dominates uniform; the
  dominance is strict whenever $\{\sigma_t\}$ are not all equal
  and $\phi_K$ is strictly concave.
\end{enumerate}
\end{corollary}

\begin{proof}
The proof of \cref{prop:offline} (\cref{sec:proof-offline}) relies on
three properties of the gain function $a_K$:
(a)~marginal gains $a_{K+1}-a_K$ are non-increasing,
(b)~$a_K$ is $t$-independent,
and (c)~$a_K' > 0$.
By \cref{thm:general-gain}, $\phi_K^{\mathcal{L}}$ inherits all
three: non-accelerating marginals from~(A2), $t$-independence
from~(A3) and~(A4), and $(\phi_K^{\mathcal{L}})'>0$ from~(A1).
(When $\phi_K$ is linear, as for local perturbation, the
allocation problem degenerates to a linear program assigning
all budget to $\arg\max_t\sigma_t$; this is a valid extreme
case of water-filling.)
Substituting $\phi_K^{\mathcal{L}}$ for $a_K$ throughout the proof
yields the result.
\end{proof}

\begin{corollary}[Online Jensen gap and dual-threshold stopping for general operators]
\label{cor:online-general}
Under the conditions of \cref{thm:general-gain}, with
prompt-dependent $\sigma_t=\sigma_t(c,x_t)$,
define the generalized oracle
\begin{equation}
  V_{\mathcal{L}}^*(\boldsymbol{\sigma})
  \;=\;
  \max_{\{K_t\}:\,\sum_t K_t=B,\;K_t\ge 1}\;
  \sum_{t=1}^{T}\sigma_t\,\phi_{K_t}^{\mathcal{L}}.
  \label{eq:oracle-general}
\end{equation}
Then:
\begin{enumerate}
\item[\textup{(i)}]
  $V_{\mathcal{L}}^*$ is convex in $\boldsymbol{\sigma}$, and
  \begin{equation}
    \mathbb{E}_{\boldsymbol{\sigma}}\!
    \bigl[V_{\mathcal{L}}^*(\boldsymbol{\sigma})\bigr]
    \;\ge\;
    V_{\mathcal{L}}^*\!\bigl(\bar{\boldsymbol{\sigma}}\bigr),
    \label{eq:jensen-general}
  \end{equation}
  with equality iff $\boldsymbol{\sigma}$ is constant across instances.
\item[\textup{(ii)}]
  The marginal gain factorizes as
  \begin{equation}
    \frac{\partial}{\partial K}\,G_t^{\mathcal{L}}(K)
    \;=\;
    \sigma_t\cdot(\phi_K^{\mathcal{L}})'.
    \label{eq:marginal-general}
  \end{equation}
  The dual-threshold stopping rule of \cref{alg:offline_online_budget}
  remains valid: stop when both $\sigma_t$ is small (sensitivity low)
  and $(\phi_K^{\mathcal{L}})'$ is small (search exhausted).
\end{enumerate}
\end{corollary}

\begin{proof}
For part~(i): for each fixed allocation $\{K_t\}$, the map
$\boldsymbol{\sigma}\mapsto\sum_t\sigma_t\,\phi_{K_t}^{\mathcal{L}}$
is linear.
$V_{\mathcal{L}}^*$ is the pointwise maximum of finitely many
linear functions, hence convex.
Jensen's inequality gives~\eqref{eq:jensen-general}.
This argument is identical to \cref{sec:proof-online} and uses
\emph{no} property of $\phi_K^{\mathcal{L}}$ beyond non-negativity.

For part~(ii): the factorization~\eqref{eq:marginal-general} follows
directly from~\eqref{eq:general-gain-factorization}.
The dual-threshold logic of \cref{prop:online}(ii) relies solely on
the product structure of the marginal gain and the
non-accelerating property of $\phi_K^{\mathcal{L}}$ (which
guarantees that the chord $\bar{g}_s\ge G_s'(\hat{K}_s)$);
both hold under (A1)--(A4).
\end{proof}

\begin{remark}[Jensen gap magnitude and operator efficiency]
\label{rmk:jensen-curvature}
The Jensen gap~\eqref{eq:jensen-general} inherits the curvature of
$\phi_K^{\mathcal{L}}$.
Faster-growing $\phi_K$ (e.g., local perturbation with linear growth)
sharpens the switching between allocation regimes as
$\boldsymbol{\sigma}$ varies, enlarging the gap:
\emph{online adaptation is more valuable for more efficient
operators}.
For pure random search ($\phi_K=a_K\sim\sqrt{2\log K}$), the slow
growth yields a smaller gap and less benefit from online
scheduling.
\end{remark}

% -----------------------------------------------------------
\subsubsection{Operator-Specific Gain Sequences}
\label{sec:operator-gain}

We now derive $\phi_K^{\mathcal{L}}$ for each practical operator and
verify that (A1)--(A4) hold.
In each case, ``linearized regime'' refers to
$g_t\sqrt{h}\to 0$, where the score function is effectively
linear in the noise: $\Psi_t(\bar{X}+g_t\sqrt{h}\,\xi;c)
= \mu_t+\sigma_t\langle u,\xi\rangle + O(g_t^2 h)$
(by \cref{prop:taylor}).

\begin{example}[Random search]
\label{ex:random-search}
The operator draws $K$ i.i.d.\ $\xi_k\sim\mathcal{N}(0,I)$
and keeps the best.
In the linearized regime, the projected scores
$\langle u,\xi_k\rangle$ are i.i.d.\ $\mathcal{N}(0,1)$, so
\begin{equation}
  \phi_K^{\mathrm{RS}} = a_K
  := \mathbb{E}\bigl[\max(Z_1,\ldots,Z_K)\bigr],
  \qquad Z_k\stackrel{\mathrm{iid}}{\sim}\mathcal{N}(0,1).
  \label{eq:phi-rs}
\end{equation}
\textbf{(A1):}~The running maximum is strictly increasing for $K\ge 2$
(a fresh draw has positive probability of exceeding the current max).
\textbf{(A2):}~$a_K$ is concave
(\citealt{david2003order}, Ch.~4).
\textbf{(A3):}~By \cref{thm:loc-scale}(iii),
$G_t(K)=\sigma_t\,a_K+O_K(g_t^2 h)$; the gain sequence
$a_K\sim\sqrt{2\log K}$ is $t$-independent.
\textbf{(A4):}~i.i.d.\ $\mathcal{N}(0,I)$ draws are
rotationally invariant.
\end{example}

\begin{example}[$\epsilon$-greedy search]
\label{ex:eps-greedy}
Starting from $\xi^{(0)}\sim\mathcal{N}(0,I)$, each iteration
draws $\xi_{\mathrm{new}}\sim\mathcal{N}(0,I)$ with probability
$\epsilon$ (explore), or sets
$\xi_{\mathrm{new}}=\xi^*+\eta\zeta$,
$\zeta\sim\mathcal{N}(0,I)$ (exploit by perturbing the current best
$\xi^*$), and updates $\xi^*$ if the new candidate is better.
In the linearized regime, the objective
$f(\xi)=\langle u,\xi\rangle$ is linear, and:
\begin{itemize}
  \item \emph{Exploration steps} ($\approx\epsilon K$ in expectation):
    each is an i.i.d.\ draw from $\mathcal{N}(0,1)$ in the projected
    coordinate $\langle u,\xi\rangle$, contributing a gain of
    $a_{\lceil\epsilon K\rceil}$.
  \item \emph{Exploitation steps} ($\approx(1{-}\epsilon)K$):
    each perturbs the current best in the direction $u$ by
    $\eta\langle u,\zeta\rangle\sim\mathcal{N}(0,\eta^2)$.
    On a linear objective, this is a random walk with positive drift
    (accept-reject filter), contributing additional gain
    beyond the exploration component.
\end{itemize}
Therefore
\begin{equation}
  \phi_K^{(\epsilon)}
  \;\ge\; a_{\lceil\epsilon K\rceil},
  \label{eq:phi-epsgreedy-lb}
\end{equation}
with equality when exploitation steps add no improvement
(e.g., when $\eta\to 0$).
In the other extreme, when $\epsilon=1$ the operator
reduces to random search and $\phi_K^{(1)}=a_K$.

\textbf{(A1):}~Each iteration has probability $\epsilon>0$ of being
an exploration step that draws a fresh
$Z_{\mathrm{new}}\sim\mathcal{N}(0,1)$ in the projected coordinate.
Because $\mathcal{N}(0,1)$ has unbounded support,
$P(Z_{\mathrm{new}}>\max_{0\le j\le K}\langle u,\xi^{(j)}\rangle)>0$,
so $\phi_{K+1}^{(\epsilon)}>\phi_K^{(\epsilon)}$ for all~$K\ge 0$.

\textbf{(A2):}~Let $M_K=\max_{0\le j\le K}\langle u,\xi^{(j)}\rangle$
denote the running best after $K$ iterations.
The marginal gain of the $(K{+}1)$-th iteration is
\[
  \Delta_K
  := \phi_{K+1}^{(\epsilon)}-\phi_K^{(\epsilon)}
  = \mathbb{E}\bigl[\max(0,\,\langle u,\xi^{(K+1)}\rangle - M_K)\bigr].
\]
With probability $\epsilon$ this iteration explores
(draws $Z_{\mathrm{new}}\sim\mathcal{N}(0,1)$ independent of $M_K$),
and with probability $1{-}\epsilon$ it exploits
(perturbs the incumbent by
$\eta\langle u,\zeta\rangle\sim\mathcal{N}(0,\eta^2)$).
In both cases, the candidate is accepted only when it exceeds~$M_K$.
For exploration, the candidate $Z_{\mathrm{new}}\sim\mathcal{N}(0,1)$
is independent of history, so
$\mathbb{E}[\max(0,Z_{\mathrm{new}}-M_K)]$ is non-increasing
in~$M_K$ (and hence in~$K$), because the tail probability
$P(Z_{\mathrm{new}}>m)$ decreases in~$m$.
For exploitation, the overshoot $\max(0,\eta\langle u,\zeta\rangle)$
does not depend on~$M_K$ at all (the perturbation is added to
the incumbent and accepted only if positive), yielding a
constant marginal $\eta\,\mathbb{E}[\max(0,Z)]=\eta/\sqrt{2\pi}$.
Therefore
$\Delta_K = \epsilon\cdot(\text{non-increasing in }K)
+ (1{-}\epsilon)\cdot(\text{constant})$,
which is non-increasing in~$K$,
so $\phi_K^{(\epsilon)}$ satisfies~(A2).

\textbf{(A3):}~In the linearized regime the projected problem
reduces to optimizing $\langle u,\xi\rangle$.
Both the exploration draw $\langle u,\xi_{\mathrm{new}}\rangle
\sim\mathcal{N}(0,1)$ and the exploitation perturbation
$\eta\langle u,\zeta\rangle\sim\mathcal{N}(0,\eta^2)$ depend
on the direction~$u$ only through a one-dimensional projection,
whose distribution is $u$-independent by isotropy of
$\mathcal{N}(0,I)$.
Hence $\phi_K^{(\epsilon)}$ does not depend on the timestep~$t$.

\textbf{(A4):}~Both primitives ($\mathcal{N}(0,I)$ draws for explore,
$\mathcal{N}(0,I)$ perturbations for exploit) are
rotationally invariant.
The accept-reject comparison uses
$\langle u,\xi\rangle$, which satisfies
$\langle Qu,Q\xi\rangle=\langle u,\xi\rangle$ for any
orthogonal~$Q$, so the full operator is rotationally equivariant.
\end{example}

\begin{example}[Local perturbation search ($\epsilon$-greedy with $\epsilon=0$)]
\label{ex:local-perturbation}
This operator is the pure-exploitation limit of the $\epsilon$-greedy
operator in \cref{ex:eps-greedy}: each iteration perturbs the
current best $\xi^*$ by a random direction
$\zeta\sim\mathcal{N}(0,I)$, setting
$\xi_{\mathrm{new}}=\xi^*+\eta\zeta$, and updates $\xi^*$ if the
new candidate scores higher.
In the linearized regime, the projected coordinate of the
incumbent evolves by accept-reject: the perturbation
$\eta\langle u,\zeta\rangle\sim\mathcal{N}(0,\eta^2)$ is
accepted when $\langle u,\xi_{\mathrm{new}}\rangle
>\langle u,\xi^*\rangle$, i.e., when $\langle u,\zeta\rangle>0$.
Because the acceptance event is independent of the current
value of the incumbent, each iteration adds a constant expected
improvement
\begin{equation}
  \phi_{K+1}^{(\mathrm{LP})}-\phi_K^{(\mathrm{LP})}
  \;=\;
  \eta\,\mathbb{E}\bigl[\max(0,Z)\bigr]
  \;=\;
  \frac{\eta}{\sqrt{2\pi}},
  \qquad Z\sim\mathcal{N}(0,1),
  \label{eq:phi-lp-marginal}
\end{equation}
giving
\begin{equation}
  \phi_K^{(\mathrm{LP})}
  \;=\;
  \frac{\eta\,K}{\sqrt{2\pi}}.
  \label{eq:phi-lp-exact}
\end{equation}

\textbf{(A1):}~Each perturbation has probability
$P(\langle u,\zeta\rangle>0)=\tfrac{1}{2}>0$ of improving the
incumbent, so $\phi_K$ is strictly increasing.

\textbf{(A2):}~In the linearized regime the marginal
gain~\eqref{eq:phi-lp-marginal} is constant (hence bounded
and non-increasing), satisfying~(A2).
In the full problem, the quadratic remainder
$R_t(\xi)=O(g_t^2 h\|\xi\|^2)$ penalizes large incumbent norms,
reducing the acceptance probability as $K$ grows and making the
marginal strictly decreasing.

\textbf{(A3):}~The perturbation
$\eta\langle u,\zeta\rangle\sim\mathcal{N}(0,\eta^2)$ is a
one-dimensional projection of $\mathcal{N}(0,I)$, whose
distribution does not depend on~$u$ or the timestep~$t$.
Hence $\phi_K^{(\mathrm{LP})}$ is $t$-independent.

\textbf{(A4):}~The perturbation $\zeta\sim\mathcal{N}(0,I)$ is
rotationally invariant, and the accept-reject comparison depends
on $\langle u,\xi\rangle$ which satisfies
$\langle Qu,Q\xi\rangle=\langle u,\xi\rangle$ for any
orthogonal~$Q$.
\end{example}

\begin{table}[t]
\centering
\caption{Summary of operator-specific gain sequences
$\phi_K^{\mathcal{L}}$ in the linearized regime.
All operators satisfy (A1)--(A4).
The growth rate determines how quickly search saturates;
faster-growing operators benefit more from online scheduling
(\cref{rmk:jensen-curvature}).}
\label{tab:phi-summary}
\begin{tabular}{@{}llll@{}}
\toprule
Operator $\mathcal{L}$ & $\phi_K^{\mathcal{L}}$ (linearized) &
Growth rate & Regime note \\
\midrule
Random search &
$a_K\sim\sqrt{2\log K}$ &
$O(\sqrt{\log K})$ &
Exact \\
$\epsilon$-greedy &
$\ge a_{\lceil\epsilon K\rceil}$ &
$\Omega(\sqrt{\log K})$ &
Lower bound \\
Local perturbation ($\epsilon{=}0$) &
$\eta K/\sqrt{2\pi}$ &
$O(K)$ &
Exact (linearized) \\
\bottomrule
\end{tabular}
\end{table}

% -----------------------------------------------------------
\subsubsection{Operator Comparison Under Fixed Budget}
\label{sec:operator-comparison}

The unified framework permits a direct comparison of operators under
a common budget.

\begin{proposition}[Crossover between random search and local perturbation]
\label{prop:crossover}
In the linearized regime, random search achieves
$\phi_K^{\mathrm{RS}}\sim\sqrt{2\log K}$,
while local perturbation achieves
$\phi_K^{(\mathrm{LP})}= \eta K/\sqrt{2\pi}$
(\cref{eq:phi-lp-exact}).
There exists a crossover budget
\begin{equation}
  K^*\;=\;\tilde{\Theta}\!\bigl(\sqrt{2\pi}/\eta\bigr)
  \label{eq:crossover}
\end{equation}
(where $\tilde{\Theta}$ suppresses logarithmic factors)
such that:
\begin{enumerate}
\item[\textup{(i)}]
  For $K\ll K^*$:
  $\phi_K^{(\mathrm{LP})}=\eta K/\sqrt{2\pi}>\sqrt{2\log K}
  =\phi_K^{\mathrm{RS}}$,
  so local perturbation dominates.
\item[\textup{(ii)}]
  For $K\gg K^*$:
  both gains continue to grow, but the quadratic remainder
  $R_t(\xi)=O(g_t^2 h\|\xi^{(K)}\|^2)$ in the full (non-linearized)
  problem increasingly penalizes local perturbation's drifted
  iterates, so the effective gains of both operators are bounded.
\end{enumerate}
\end{proposition}

\begin{proof}
The crossover occurs at the unique solution of
$\eta K/\sqrt{2\pi}=\sqrt{2\log K}$.
Since $\eta K/\sqrt{2\pi}$ grows linearly and $\sqrt{2\log K}$
grows logarithmically, there exists a unique crossing point.
Substituting $K=\sqrt{2\pi}\sqrt{2\log K}/\eta$ and
iterating once gives
$K^*=\sqrt{2\pi}\sqrt{2\log(\sqrt{2\pi}/\eta)}/\eta
\cdot(1+o(1))$ as $1/\eta\to\infty$, which is
$\tilde{\Theta}(\sqrt{2\pi}/\eta)$.
Below this crossover, the linear growth $\eta K/\sqrt{2\pi}$
of local perturbation exceeds the logarithmic growth
$\sqrt{2\log K}$ of random search.
Above it, local perturbation continues to grow linearly in the
linearized regime, but the quadratic remainder
$R_t(\xi)=O(g_t^2 h\|\xi^{(K)}\|^2)$ in the full problem
increasingly penalizes the drifted iterates, effectively
bounding the achievable gain for both operators.
\end{proof}

\begin{remark}[Operator selection as part of scheduling]
\label{rmk:operator-selection}
\cref{prop:crossover} implies that the optimal operator choice
depends on the per-step budget $K_t$: high-value timesteps receiving
concentrated budget benefit from local perturbation's linear
initial growth,
while low-value timesteps see negligible difference between operators.
Extending the scheduling decision to include operator selection is a
natural next step; the MDP formulation of \cref{sec:mdp} already
accommodates this as a richer action space.
\end{remark}
